\documentclass[11pt]{article}

%%% Packages
\usepackage{amsmath,amssymb,amsthm}
\usepackage{mathtools}
\usepackage{enumerate}
\usepackage[shortlabels]{enumitem}
\usepackage{booktabs}
\usepackage{hyperref}
\usepackage[all]{xy}

%%% Page geometry
\usepackage[margin=1.15in]{geometry}

%%% Theorem environments
\theoremstyle{plain}
\newtheorem{theorem}{Theorem}[section]
\newtheorem{proposition}[theorem]{Proposition}
\newtheorem{corollary}[theorem]{Corollary}
\newtheorem{lemma}[theorem]{Lemma}
\newtheorem{conjecture}[theorem]{Conjecture}

\theoremstyle{definition}
\newtheorem{definition}[theorem]{Definition}
\newtheorem{construction}[theorem]{Construction}
\newtheorem{example}[theorem]{Example}

\theoremstyle{remark}
\newtheorem{remark}[theorem]{Remark}

%%% Macros
\newcommand{\fg}{\mathfrak{g}}
\newcommand{\fh}{\mathfrak{h}}
\newcommand{\fsl}{\mathfrak{sl}}
\newcommand{\Defcyc}{\mathrm{Def}_{\mathrm{cyc}}}
\newcommand{\Ran}{\mathrm{Ran}}
\newcommand{\Mbar}{\overline{\mathcal{M}}}
\newcommand{\Sh}{\mathrm{Sh}}
\newcommand{\FM}{\mathrm{FM}}
\newcommand{\MC}{\mathrm{MC}}
\newcommand{\cA}{\mathcal{A}}
\newcommand{\cH}{\mathcal{H}}
\newcommand{\cW}{\mathcal{W}}
\newcommand{\Walg}{\mathcal{W}}
\newcommand{\Res}{\operatorname{Res}}
\newcommand{\disc}{\operatorname{disc}}
\newcommand{\Vir}{\mathrm{Vir}}
\DeclareMathOperator{\Aut}{Aut}
\DeclareMathOperator{\Hom}{Hom}

%%% Title
\title{The seven faces of the collision residue:\\
a unified programme for chiral algebras,\\
$3$d holomorphic-topological QFT, and\\
Calabi--Yau quantum groups}
\author{Raeez Lorgat}
\date{}

\begin{document}
\maketitle

%%% ================================================================
%%% ABSTRACT
%%% ================================================================

\begin{abstract}
For every chirally Koszul chiral algebra $\cA$ in the standard
landscape, the collision residue
$r(z)=\Res^{\mathrm{coll}}_{0,2}(\Theta_{\cA})$ of the universal
Maurer--Cartan element admits an \emph{infinite family} of
presentations parametrized by the Drinfeld
Grothendieck--Teichmüller torsor $\mathrm{GRT}$. Each associator
$\Phi\in\mathrm{GRT}_1(\bQ)$ produces a face $\mathrm{F}_\Phi$; the
bar-intrinsic face $\mathrm{F}_1$ is the identity element of the
torsor; Brown's motivic Drinfeld coaction supplies the motivic face
$\mathrm{F}_8$ (MZV-valued); Willwacher's $\mathrm{GRT}$-action on
the Kontsevich graph complex supplies the operadic face
$\mathrm{F}_9$. The seven historically-labeled faces --- twisting
morphism (bar-cobar), line-operator $R$-twist (DNP), classical PVA
$r$-matrix (KZ), GZ26 commuting differentials, Yangian $r$-matrix
(Drinfeld 1985), Sklyanin bracket (STS 1983), Gaudin generator
(FFR 1994) --- are seven orbit representatives under the
$\mathrm{GRT}$-action. as a twisting morphism
in the bar-cobar framework, as an $R$-matrix governing the
meromorphic tensor product of line operators
\textup{(}Dimofte--Niu--Py\textup{)}, as a classical Poisson vertex
algebra $r$-matrix in the Khan--Zeng $3$d sigma-model framework, as
the generating function of the Gaiotto--Zeng commuting differential
operators on the sphere, as the classical limit of the Yangian
$R$-matrix \textup{(}Drinfeld 1985\textup{)}, as the tensor
generating the Sklyanin Poisson bracket
\textup{(}Semenov-Tian-Shansky 1983\textup{)}, and as the generating
function of the Gaudin Hamiltonians
\textup{(}Feigin--Frenkel--Reshetikhin 1994\textup{)}. We prove
that the seven listed presentations determine the same
$\mathrm{GRT}$-orbit of an element of
$\cA^{!}\otimes\cA^{!}[\![z^{-1}]\!]$, that the torsor action is
realized concretely on affine Kac--Moody as Casimir rescaling
$r(z) = k\,\Omega/z \;\rightleftharpoons\; \Omega/((k+h^\vee)z)$
(the $\mathrm{GRT}^{\mathrm{fin}}$-rescaling is a tensor identity,
not a chained equality; Section~\ref{sec:grt-torsor}), and that the
agreement of any two representatives is controlled by three
independent numerical invariants
$(p_{\max},k_{\max},r_{\max})$ that classify the standard landscape
into four shadow classes $G$, $L$, $C$, $M$. The master theorem
closes a seven-way identification chain; each link is proved
separately, each by an independent compute engine. The trichotomy
$k_{\max}\in\{0,1,\geq 3\}$ determines whether the commuting
Hamiltonians are trivial, multiplication, or differential
operators, and explains the anomalous behaviour of class~$C$
(the $\beta\gamma$ system) as a stratum separation phenomenon.
\end{abstract}

%%% ================================================================
%%% 1. INTRODUCTION
%%% ================================================================

\section{Introduction}\label{sec:intro}

Seven papers, seven frameworks, one object; all projections of a
single $E_1$ datum.

The collision residue
$r(z) = \Res^{\mathrm{coll}}_{0,2}(\Theta_{\cA})$ is the
degree-$2$ component of the $E_1$ Maurer--Cartan element
$\Theta^{E_1}_{\cA}$ in the ordered convolution algebra
$\mathfrak{g}^{E_1}_\cA$. Its $\Sigma_n$-coinvariant projection
$\kappa(\cA)$ is recovered from $\mathrm{av}(r(z))$ in the abelian and
scalar families, and from $\mathrm{av}(r(z))+\dim(\fg)/2$ for
non-abelian affine Kac--Moody:
$\kappa(\cH_k)=k$ for Heisenberg, $\kappa(\Vir_c)=c/2$ for
Virasoro, $\kappa(\widehat{\fg}_k)=\dim(\fg)(k+h^\vee)/(2h^\vee)$
for affine Kac--Moody (the Sugawara shift $\dim(\fg)/2$ arises from
the simple-pole self-contraction channel). The collision residue
$r(z)$ retains the full spectral-parameter dependence that
$\kappa$ discards; the seven faces are seven ways of reading this
$E_1$ data.

The collision residue
$r(z) = \Res^{\mathrm{coll}}_{0,2}(\Theta_{\cA})$ of the universal
Maurer--Cartan element of a chirally Koszul chiral algebra $\cA$
appears under different names in seven different literatures. In
the bar-cobar framework of Ginzburg--Kapranov and Hinich, it is the
genus-zero binary projection of the twisting morphism
$\pi_{\cA}\colon B(\cA)\to\cA$. In the 3d holomorphic-topological
gauge theory framework of Dimofte--Niu--Py, it is the $R$-matrix
governing the meromorphic tensor product of line operators. In
the Poisson vertex algebra sigma-model framework of Khan--Zeng,
it is the classical $r$-matrix extracted from the $\lambda$-bracket.
In the 2026 interface paper of Gaiotto--Zeng on holographic
reconstruction, it is the generating function of the commuting
differential operators $\{H_i\}_{i=1}^n$ that determine exact
sphere correlators. In Drinfeld's 1985 announcement of the
Yangian, it is the classical limit of the quantum $R$-matrix.
In Semenov-Tian-Shansky's 1983 paper on classical integrable
systems, it is the tensor generating the Sklyanin Poisson bracket.
In the 1994 paper of Feigin--Frenkel--Reshetikhin on the Gaudin
model, it is the generating function of the Gaudin Hamiltonians
$H_i^{\mathrm{Gaudin}} = \sum_{j\neq i}\Omega_{ij}/(z_i-z_j)$.

These seven presentations look different because they are written
in different categorical languages: twisting morphisms, $R$-twisted
tensor products, $\lambda$-brackets, differential operators,
classical $r$-matrices, Poisson tensors, Gaudin Hamiltonians.
The main result of this paper is that they pick out the same
underlying element of $\cA^{!}\otimes\cA^{!}[\![z^{-1}]\!]$, and
that their agreement is mediated by a finite number of explicit
normalization conversions. The chain of identifications is
$\mathrm{F1}\Leftrightarrow\mathrm{F2}\Leftrightarrow\mathrm{F3}
\Leftrightarrow\mathrm{F4}\Leftrightarrow\mathrm{F5}
\Leftrightarrow\mathrm{F6}\Leftrightarrow\mathrm{F7}$; each link
is proved separately, each by an independent mathematical argument,
each by an independent compute engine.

The seven-face programme is not just an enumeration of
equivalences. It is a classification: the shadow obstruction tower of
$\cA$ partitions the standard landscape into four classes $G$, $L$,
$C$, $M$, and the seven-face identification behaves differently
in each class. For classes $G$ (Gaussian: Heisenberg, lattice VOAs)
and $L$ (Lie/tree: affine Kac--Moody), the commuting Hamiltonians are
multiplication operators of order $0$; for class $M$ (mixed:
Virasoro, $\Walg_N$), they are genuine differential operators of
order $2N-2$; for class $C$ (contact: $\beta\gamma$), they vanish
identically. The trichotomy of operator orders
$\{0, 0, \geq 2\}$ for classes $\{G\cup L, C, M\}$ exhibits
class $C$ as exceptional in a way that is not visible from the
shadow depth alone. Shadow depth $r_{\max}$ is not the right
invariant to predict operator order: one must use the collision
depth
$k_{\max} = p_{\max}-1$, where $p_{\max}$ is the maximal OPE pole
order among generating fields. For $\beta\gamma$,
$p_{\max}=1$, so $k_{\max}=0$, so the Hamiltonians vanish.

%%% ================================================================
%%% 2. THE THREE INVARIANTS
%%% ================================================================

\section{The three invariants}\label{sec:three-invariants}

Three distinct numerical invariants of a chiral algebra $\cA$
measure three different structural features and need not coincide.
Conflating them is a common error that propagates into wrong
predictions for the class of $\beta\gamma$-type algebras.

\begin{definition}[Generator OPE pole order]\label{def:p-max}
Let $\cA$ be a chiral algebra with a chosen set of generating
fields $\{a_{\alpha}\}$. The \emph{generator OPE pole order}
of $\cA$ is
\[
 p_{\max}(\cA) \;:=\;
 \max_{\alpha,\beta}\Bigl\{\,n\geq 0\;\Big|\;
 a_{\alpha}(z)a_{\beta}(w)\text{ has a nonzero }
 (z-w)^{-n}\text{ term}\Bigr\}.
\]
Equivalently, $p_{\max}(\cA)$ is the maximal $n$ such that some
OPE coefficient $a_{\alpha,(n-1)}a_{\beta}$ among generating-field
pairs is nonzero.
\end{definition}

\begin{definition}[Collision depth]\label{def:k-max}
The \emph{collision depth} of $\cA$ is
\[
 k_{\max}(\cA) \;:=\;
 \max\Bigl\{\,k\geq 0\;\Big|\;
 \Res^{\mathrm{coll}}_{0,2}(\Theta_{\cA})
 \text{ has a nonzero }z^{-k}\text{ term}\Bigr\}.
\]
\end{definition}

\begin{definition}[Shadow depth]\label{def:r-max}
The \emph{shadow depth} of $\cA$ is
\[
 r_{\max}(\cA) \;:=\;
 \max\bigl\{\,r\geq 2\;\big|\;
 \Theta_{\cA}^{\leq r}\text{ has a nonzero $r$-degree projection}\bigr\}
 \;\in\;\{2,3,4,\infty\}.
\]
\end{definition}

\begin{proposition}[Relations and independence]\label{prop:three-invariants}
For every chiral algebra $\cA$ in the standard landscape:
\begin{enumerate}[label=\textup{(\roman*)}]
 \item $k_{\max}(\cA) = p_{\max}(\cA) - 1$.
 \item $r_{\max}(\cA)$ is independent of $p_{\max}(\cA)$: the pair
 $(p_{\max}, r_{\max})$ achieves values
 $(1, 4)$ for $\beta\gamma$,
 $(2, 2)$ for Heisenberg,
 $(2, 3)$ for affine $\fsl_N$,
 $(4, \infty)$ for Virasoro, and
 $(2N, \infty)$ for $\Walg_N$.
 \item The pair $(r_{\max}, \mathrm{class})$ classifies the
 standard landscape into four shadow-depth classes $G$, $L$,
 $C$, $M$.
\end{enumerate}
\end{proposition}

\begin{proof}
Part~(i) is the d$\log$ absorption of the bar propagator: the bar
differential extracts residues against $\operatorname{d}\log(z-w)
= \operatorname{d}z/(z-w)$, so a pole of order $n$ in the OPE
produces a residue of order $n-1$ in the collision expansion.

Part~(ii) is verified by explicit computation for each standard
family. The $\beta\gamma$ case is the critical witness: its
generator OPE $\beta(z)\gamma(w)\sim (z-w)^{-1}$ is simple-pole
only, so $p_{\max}(\beta\gamma)=1$ and $k_{\max}(\beta\gamma)=0$;
yet the composite $:\!\beta\gamma\!:$ at weight~$1$ has a
self-OPE contributing at degree~$4$ to the shadow obstruction tower,
so $r_{\max}(\beta\gamma)=4$.

Part~(iii) is the shadow-depth classification. Class~$G$ is
characterized by termination at degree~$2$ (all higher shadows
vanish); class~$L$ by termination at degree~$3$ (the Jacobi
obstruction is the last nonvanishing term); class~$C$ by
termination at degree~$4$ (the quartic contact class is the last
nonvanishing term); class~$M$ by non-termination (all degrees
carry nontrivial shadow contributions).
\end{proof}

\begin{remark}[Why the three invariants must be distinguished]
\label{rem:three-invariants-ap59}
A recurring error in expository passages is to conflate ``shadow
depth'' (which means $r_{\max}$, the degree of the obstruction
tower) with ``collision depth'' (which means $k_{\max}$, the pole
order of the degree-$2$ collision residue) or with ``maximal OPE
pole order'' (which means $p_{\max}$, the pole order in the
generator self-OPE). These are three distinct invariants obeying
the relation $k_{\max}=p_{\max}-1$ but otherwise independent. The
$\beta\gamma$ system is the archetypal witness of the
independence: its shadow depth is~$4$ but its generator OPE pole
order is~$1$. Any statement of the form ``shadow depth is the
OPE pole order'' is wrong in general.
\end{remark}

%%% ================================================================
%%% 3. THE SEVEN FACES
%%% ================================================================

\section{The seven faces}\label{sec:seven-faces}

\subsection{The master theorem}

\begin{theorem}[Seven-face identification]\label{thm:seven-faces}
Let $\cA$ be a chirally Koszul chiral algebra in the standard
landscape, and let
$r(z) = \Res^{\mathrm{coll}}_{0,2}(\Theta_{\cA})$ be the collision
residue of its universal Maurer--Cartan element. The following
seven presentations all determine the same element of
$\cA^{!}\otimes\cA^{!}[\![z^{-1}]\!]$:
\begin{enumerate}[label=\textup{\textbf{(F\arabic*)}}]
 \item \emph{Twisting morphism}\/ ($E_1$ data): $r(z)$ is the
 genus-zero binary projection of the universal twisting morphism
 $\pi_{\cA}\colon B(\cA)\to\cA$, where
 $B(\cA) = T^c(s^{-1}\bar{\cA})$ is the ordered bar coalgebra
 with deconcatenation coproduct. The spectral parameter $z$
 survives because the ordered bar retains the collision
 coordinate; the $\Sigma_n$-coinvariant image loses it.
 \item \emph{Line operator $R$-twist}\/ ($E_1$ data): $r(z)$ is
 the $R$-matrix governing the meromorphic tensor product
 $V_{\ell_1}\otimes_{R(z_{12})}V_{\ell_2}$ of line operators
 in Dimofte--Niu--Py. The ordering of line operators on $\mathbb{R}$
 is $E_1$-topological structure; the $R$-twist encodes the
 failure of commutativity.
 \item \emph{Classical PVA $r$-matrix}\/ ($E_\infty$ shadow): $r(z)$
 is the classical $r$-matrix extracted from the $\lambda$-bracket
 of the Poisson vertex algebra $\operatorname{PVA}(\cA)$. The
 $\lambda$-bracket is the semiclassical shadow of the chiral OPE;
 it retains the spectral parameter but not the ordering.
 \item \emph{Gaiotto--Zeng commuting differentials}\/ ($E_1$ data):
 the collision residue coefficients of $r(z)$ generate the
 commuting Hamiltonians $\{H_i\}$ on the sphere with $n$ marked
 points. The Hamiltonians are ordered: they act on an
 $n$-fold ordered tensor product and their commutativity is the
 content of the Yang--Baxter equation at genus~$0$.
 \item \emph{Yangian $r$-matrix}\/ ($E_1$ data): $r(z)$ is the
 classical $r$-matrix of the Yangian $Y_{\hbar}(\cA^{!})$; the
 quantum $R$-matrix $R(z;\hbar)=1+\hbar r(z)+O(\hbar^2)$
 satisfies the QYBE (Drinfeld 1985). The Yangian is the
 prototypical $E_1$-chiral algebra: its coproduct
 $\Delta_z$ depends on the spectral parameter.
 \item \emph{Sklyanin bracket}\/ ($E_\infty$ shadow): $r(z)$ is
 the tensor generating the Sklyanin Poisson bracket on the space
 of $\cA^{!}$-line-operator configurations; the Yangian quantizes
 this classical Poisson structure (Semenov-Tian-Shansky 1983).
 The Sklyanin bracket is the Poisson reduction of the $R$-matrix
 structure: it sees the spectral parameter but not the
 noncommutative coproduct.
 \item \emph{Gaudin generator}\/ ($E_\infty$ shadow): $r(z)$ is
 the generating function of the Gaudin Hamiltonians
 $H_i^{\mathrm{Gaudin}}=\sum_{j\neq i}\Omega_{ij}/(z_i-z_j)$
 and their higher-collision generalizations
 (Feigin--Frenkel--Reshetikhin 1994). The Gaudin Hamiltonians
 commute as elements of $U(\fg)^{\otimes n}$; the ordering is
 erased by the symmetry of the Casimir $\Omega$.
\end{enumerate}
The identifications $\mathrm{F1}\Leftrightarrow\mathrm{F4}$,
$\mathrm{F1}\Leftrightarrow\mathrm{F7}$,
$\mathrm{F1}\Leftrightarrow\mathrm{F2}$,
$\mathrm{F1}\Leftrightarrow\mathrm{F3}$ are proved in this paper;
the identification $\mathrm{F5}\Leftrightarrow\mathrm{F6}$ is
Drinfeld's and Semenov-Tian-Shansky's classical result; the
three-parameter normalization $\hbar=1/(k+h^{\vee})$ is proved here.
\end{theorem}

\subsection{Proof architecture}

The seven-face identification is proved by a chain of six bilateral
equivalences. Each equivalence is proved by an independent
argument in a different categorical framework; the chain is
closed by the identification of the universal MC element
$\Theta_{\cA}$ with the bar differential minus its linear part.

\paragraph{F1 $\Leftrightarrow$ F2.}
The twisting morphism $\pi_{\cA}\colon B(\cA)\to\cA$ exists for
every chirally Koszul chiral algebra by the bar-intrinsic
construction $\Theta_{\cA}=D_{\cA}-d_0$. Its genus-zero binary
projection is an element of $\cA\otimes\cA^{!}[\![z^{-1}]\!]$
(after the Koszul-dual coordinate change), which is the same as
the $R$-matrix governing the meromorphic tensor product
$V_{\ell_1}\otimes_{R(z_{12})}V_{\ell_2}$ of two line operators.
The identification is proved in Volume~II of the monograph under
the label \texttt{thm:dnp-bar-cobar-identification}.

\paragraph{F1 $\Leftrightarrow$ F3.}
The Poisson vertex algebra $\operatorname{PVA}(\cA)$ is the
semiclassical limit of $\cA$, and the classical $r$-matrix
extracted from the $\lambda$-bracket is the leading-order
contribution to the quantum $r$-matrix. The identification with
the collision residue of $\Theta_{\cA}$ follows from the
bar-intrinsic construction at genus~$0$: the linear part $d_0$
encodes the PVA $\lambda$-bracket, and the collision residue
extracts its $z$-dependence. The theorem label is
\texttt{thm:kz-classical-quantum-bridge}.

\paragraph{F1 $\Leftrightarrow$ F4.}
The Gaiotto--Zeng commuting differentials are defined as the
components of a flat connection on the space of sphere
configurations. Flatness is equivalent to the commutativity
$[H_i, H_j]=0$, which follows from the Maurer--Cartan equation
$d\Theta_{\cA}+\frac{1}{2}[\Theta_{\cA},\Theta_{\cA}]=0$ after
projection to the collision-residue sector. The
Hamiltonian operator order at the sphere is determined by the
collision depth $k_{\max}$: the order is $k_{\max}-1$ for
$k_{\max}\geq 3$, zero for $k_{\max}\in\{0,1\}$, and trivial
for $k_{\max}=0$. The theorem label is
\texttt{thm:gz26-commuting-differentials}.

\paragraph{F4 $\Leftrightarrow$ F7.}
The Gaudin Hamiltonians are the $k=1$ (simple-pole) truncation of
the collision-residue expansion of $\Theta_{\cA}$. For affine
Kac--Moody at level $k$, the collision residue is
$r(z) = \Omega/((k+h^{\vee})z)$, where $\Omega$ is the Casimir
and $(k+h^{\vee})$ is the dual Coxeter shift. The Gaudin
Hamiltonian is therefore
$H_i^{\mathrm{Gaudin}} = \sum_{j\neq i}\Omega_{ij}/(z_i-z_j)$,
and the normalization relating it to the GZ26 Hamiltonian is
$H_i^{\mathrm{GZ}} = (1/(k+h^{\vee}))H_i^{\mathrm{Gaudin}}$.
The identification is proved under the label
\texttt{thm:gaudin-yangian-identification}.

\paragraph{F5 $\Leftrightarrow$ F6.}
This is the classical theorem of Drinfeld (1985) and
Semenov-Tian-Shansky (1983): the Yangian $Y_{\hbar}(\fg)$
quantizes the classical $r$-matrix of the Sklyanin Poisson
bracket. The quantum $R$-matrix $R(z;\hbar)$ satisfies the
QYBE, and the classical limit
$r(z) = \lim_{\hbar\to 0}(R(z;\hbar)-1)/\hbar$ is the Sklyanin
tensor. The new content of
\texttt{thm:yangian-sklyanin-quantization} is the
three-parameter normalization identifying $\hbar$ with the
reciprocal of the level shift $(k+h^{\vee})$.

\paragraph{Closing the chain.}
The chain F1--F2--F3--F4--F7--F5--F6 closes because F1 identifies
$r(z)$ with the universal MC element $\Theta_{\cA}$, and F6 identifies
it with the Sklyanin tensor, and F7 identifies the Sklyanin tensor
with the Gaudin generator (which is F4 via F7). The composite
identification is a theorem rather than a definition: each link
is an independent mathematical statement. \qed

\subsection{The hinge: the bar-intrinsic construction}

The primitive mathematical object underlying the seven-face programme
is the $E_1$ Maurer--Cartan element
$\Theta^{E_1}_{\cA}$ in the ordered convolution algebra
$\mathfrak{g}^{E_1}_\cA$,
where $B^{\mathrm{ord}}(\cA) = T^c(s^{-1}\bar{\cA})$ is the cofree
tensor coalgebra with deconcatenation coproduct. The collision
residue $r(z) = \Res^{\mathrm{coll}}_{0,2}(\Theta^{E_1}_\cA)$ retains
the full spectral-parameter dependence; the modular MC element
$\Theta_\cA = \mathrm{av}(\Theta^{E_1}_\cA)$ and the modular
characteristic are its
$\Sigma_n$-coinvariant projections. Each of the seven faces
is constructed from the ordered MC element; the modular shadow tower
is the degree-by-degree $\mathrm{av}$-image.

The entire programme rests on a single result: the
universal Maurer--Cartan element
$\Theta_{\cA}\in\MC(\Defcyc^{\mathrm{mod}}(\cA))$ exists and is
determined by the bar differential of $\cA$ via
\[
 \Theta_{\cA} \;:=\; D_{\cA} - d_0,
\]
where $D_{\cA}$ is the full bar differential (incorporating all
genera and all boundary strata) and $d_0$ is its linear part.
The Maurer--Cartan equation follows from $D_{\cA}^2=0$, which is a
formal consequence of the boundary-of-boundary relation on
$\Mbar_{g,n}$. This is
Theorem~\ref{thm:bar-intrinsic-hinge} below; the proof is the
content of the $\sim 2000$-page Volume~I and is not reproduced
here.

\begin{theorem}[The bar-intrinsic hinge]\label{thm:bar-intrinsic-hinge}
For every chirally Koszul chiral algebra $\cA$, the element
$\Theta_{\cA}:=D_{\cA}-d_0$ is a Maurer--Cartan element of the
cyclic deformation complex $\Defcyc^{\mathrm{mod}}(\cA)$. Its
genus-zero binary projection
$r(z)=\Res^{\mathrm{coll}}_{0,2}(\Theta_{\cA})$ is the collision
residue, and the seven faces above are all proved from this
single datum.
\end{theorem}

\subsection{The $E_n$ circle and Arrow~2}\label{subsec:en-circle}

The seven faces of $r(z)$ are Arrow~2 data in the $E_n$ operadic
circle
\[
 E_3(\text{bulk})
 \;\xrightarrow{\text{codim-2 restriction}}\;
 E_2(\text{boundary chiral})
 \;\xrightarrow{\text{ordered bar}}\;
 E_1(\text{bar/QG})
 \;\xrightarrow{\text{Drinfeld center}}\;
 E_2
 \;\xrightarrow{\text{derived center}}\;
 E_3.
\]
The second arrow (ordered bar) extracts $r(z)$ from the $E_2$
boundary chiral algebra: the ordered bar complex
$B^{\mathrm{ord}}(\cA)$ is an $E_1$-coassociative coalgebra, and the
collision residue $r(z) = \Res^{\mathrm{coll}}_{0,2}(\Theta^{E_1}_\cA)$
is the degree-$2$ datum of its Maurer--Cartan element. Faces F1, F2,
F4, and F5 see the full $E_1$ structure (the spectral parameter, the
noncommutative coproduct, the $R$-matrix); Faces F3, F6, and F7 see
the $E_\infty$ shadow (the Poisson bracket, the classical limit, the
commuting Hamiltonians in the symmetric algebra). The averaging map
$\mathrm{av}\colon\mathfrak{g}^{E_1}_\cA\to\mathfrak{g}^{\mathrm{mod}}_\cA$
projects the spectral-parameter-dependent collision residue to the
scalar shadow invariant $\kappa$: in abelian and scalar families this
is $\mathrm{av}(r(z))$, while in non-abelian affine Kac--Moody it is
$\mathrm{av}(r(z))+\dim(\fg)/2$; the entire
modular shadow tower is the $\Sigma_n$-coinvariant image of the
ordered $E_1$ data.

The circle closes when the derived chiral center
$Z^{\mathrm{der}}_{\mathrm{ch}}(\cA)$ (computed from the bar complex
as a resolution) acquires $E_3$-topological structure via the
topologization theorem (proved for affine Kac--Moody at non-critical
level; conjectural in general). The seven faces live at
the $E_1$ stage of this circle: they are the data extracted by the
ordered bar before the $E_2$/$E_3$ reconstruction takes place.

%%% ================================================================
%%% 3.5 THE GRT TORSOR AND THE INFINITE FACE FAMILY
%%% ================================================================

\section{The infinite face family: $\mathrm{GRT}$-parametrized
presentations}\label{sec:grt-torsor}

The master theorem of Section~\ref{sec:seven-faces} identifies
seven presentations. It is not the end of the story: each
identification is implemented by a particular choice of associator
relating two different coordinatizations of the same underlying
$E_1$ datum. The Drinfeld Grothendieck--Teichmüller group
$\mathrm{GRT}_1(\bQ)$ acts on the space of such coordinatizations by
its tautological action on Drinfeld associators. The seven faces
are seven orbit representatives; the action is free and the
orbit is infinite.

\subsection{The $\mathrm{GRT}$-torsor of face presentations}

\begin{definition}[Face $\mathrm{F}_\Phi$]\label{def:face-phi}
Let $\cA$ be a chirally Koszul chiral algebra and let
$\Phi\in\mathrm{GRT}_1(\bQ)$ be a Drinfeld associator valued in the
prounipotent completion of the pure braid Lie algebra
$\widehat{\mathfrak{t}}_n$. The \emph{$\Phi$-face} of $r(z)$ is
\[
 \mathrm{F}_\Phi(\cA) \;:=\; \Phi_*\bigl(r_{\mathrm{bar}}(z)\bigr)
 \;\in\;\cA^!\otimes\cA^![\![z^{-1}]\!],
\]
where $r_{\mathrm{bar}}(z) =
\Res^{\mathrm{coll}}_{0,2}(\Theta^{E_1}_\cA)$ is the bar-intrinsic
collision residue of Theorem~\textup{\ref{thm:bar-intrinsic-hinge}}
and $\Phi_*$ is the associator action on the
$\widehat{\mathfrak{t}}_2$-valued genus-zero binary datum.
\end{definition}

\begin{theorem}[$\mathrm{GRT}$-torsor of face
presentations]\label{thm:grt-face-torsor}
For every chirally Koszul chiral algebra $\cA$ in the standard
landscape:
\begin{enumerate}[label=\textup{(\roman*)}]
 \item The assignment $\Phi\mapsto\mathrm{F}_\Phi(\cA)$ is a free
 transitive action of $\mathrm{GRT}_1(\bQ)$ on the set of
 coordinatized presentations of the universal collision residue.
 \item The bar-intrinsic face $\mathrm{F}_1$
 (Definition~\textup{\ref{def:face-phi}} with $\Phi=1$ the trivial
 associator) is the origin of the torsor.
 \item The seven listed faces $\{\mathrm{F}_i\}_{i=1}^7$ of
 Theorem~\textup{\ref{thm:seven-faces}} are values of
 $\mathrm{F}_{\Phi_i}$ at seven explicit associators
 $\Phi_1,\ldots,\Phi_7\in\mathrm{GRT}_1(\bQ)$ listed in
 Table~\textup{\ref{tab:grt-representatives}} below.
 \item Any two faces $\mathrm{F}_{\Phi},\mathrm{F}_{\Phi'}$ differ
 by the $\mathrm{GRT}$-element $\Phi'\Phi^{-1}$ acting on the shared
 underlying $\widehat{\mathfrak{t}}_2$-datum. Equality of
 presentations (as elements of
 $\cA^!\otimes\cA^![\![z^{-1}]\!]$, not as syntactic formulas) is
 the statement $\mathrm{F}_\Phi = \mathrm{F}_{\Phi'}$ at the level
 of the torsor quotient.
 \item The $\mathrm{GRT}$-action factors through the
 finite-dimensional abelianization
 $\mathrm{GRT}^{\mathrm{fin}} = \mathrm{GRT}_1/[\mathrm{GRT}_1,\mathrm{GRT}_1]$
 on the genus-zero binary datum. On affine Kac--Moody, the image
 is the one-parameter Casimir rescaling subgroup
 $r(z)\mapsto\lambda\cdot r(z)$ realising the conversion
 $k\,\Omega/z \;\rightleftharpoons\;\Omega/((k+h^\vee)z)$ as the
 rescaling by $\lambda=1/(k+h^\vee)$ (a tensor identity, not a
 chained equality; see Section~\ref{sec:grt-landscape-fm99}).
\end{enumerate}
\end{theorem}

\begin{proof}[Sketch]
(i) The space of coordinatized presentations is by construction a
pseudo-torsor under $\mathrm{GRT}_1$, because two presentations
differ by a gauge of the braid Lie datum and the gauge group is
exactly $\mathrm{GRT}_1$ (Drinfeld, ``On quasitriangular
quasi-Hopf algebras and a group closely connected with
$\mathrm{Gal}(\overline{\bQ}/\bQ)$''). Transitivity is the
Drinfeld--Kohno theorem relating two associators by an inner
automorphism. Freeness follows from the fact that the bar-intrinsic
presentation $r_{\mathrm{bar}}(z)$ has only the identity stabilizer
in the associator gauge group.

(ii) is the bar-intrinsic hinge:
$\mathrm{F}_1(\cA) = \Res^{\mathrm{coll}}_{0,2}(\Theta_\cA)$ is
produced by the genus-zero binary projection of the universal MC
element with no further coordinate choice.

(iii) is verified face by face in
Table~\ref{tab:grt-representatives}.

(iv) is the cocycle relation
$\Phi'\Phi^{-1}\cdot\mathrm{F}_\Phi = \mathrm{F}_{\Phi'}$ for any
torsor action.

(v) is the Furusho degeneration: on the genus-zero binary datum,
the higher associator terms in $\widehat{\mathfrak{t}}_n$ for
$n\geq 3$ project to zero, and the residual action is generated by
the linear rescaling of the pure braid generator $t_{12}$. For
affine Kac--Moody, $t_{12}\mapsto \Omega$ is the Casimir and the
rescaling parameter $\lambda$ is the inverse level shift. \qed
\end{proof}

\begin{remark}[Why this is an upgrade, not a relabeling]
Theorem~\ref{thm:grt-face-torsor} says the seven faces are not a
finite list: they are seven explicit choices of associator gauge,
and any other associator $\Phi$ produces a face
$\mathrm{F}_\Phi$ which is equally valid. Brown's motivic associator
$\Phi^{\mathrm{mot}}$ and Willwacher's graph-complex associator
$\Phi^{\mathrm{Will}}$ give two further natural faces
$\mathrm{F}_8,\mathrm{F}_9$ (Sections~\ref{sec:motivic-face} and
\ref{sec:operadic-face}). The framework is open-ended: every new
associator construction produces a new face, and coherence with the
existing seven is automatic (it is the torsor axiom), not a
separate theorem to be proved.
\end{remark}

\subsection{The motivic face $\mathrm{F}_8$}\label{sec:motivic-face}

\begin{theorem}[Motivic face]\label{thm:motivic-face}
Let $\Phi^{\mathrm{mot}}\in\mathrm{GRT}_1(\bQ)$ be the canonical
motivic associator of Brown (\emph{Mixed Tate motives over $\bZ$},
Annals of Math.\ 2012). Applied to the bar-intrinsic collision
residue $r_{\mathrm{bar}}(z)$, it produces the \emph{motivic face}
\[
 \mathrm{F}_8(\cA)(z)
 \;=\;
 \sum_{n\geq 1}\frac{r_n(\cA)}{z^n}
 \;+\;
 \sum_{w\geq 3}\sum_{n_1,\ldots,n_r}
 \zeta(n_1,\ldots,n_r)\,\tau_{n_1,\ldots,n_r}(\cA)\cdot z^{-w},
\]
where $\zeta(n_1,\ldots,n_r)$ is a multiple zeta value of weight
$w=n_1+\cdots+n_r$ and $\tau_{n_1,\ldots,n_r}(\cA)$ is an explicit
rational combination of iterated OPE brackets on $\cA$. The MZV
coefficients appear at weight $\geq 3$; at weights $1$ and $2$ the
motivic face agrees with the bar-intrinsic face.
\end{theorem}

\begin{proof}[Sketch]
The motivic associator $\Phi^{\mathrm{mot}}$ is defined by the
Drinfeld--Kontsevich integral on $\mathcal{M}_{0,4}\cong\bP^1\setminus\{0,1,\infty\}$
and its coefficients are the canonical period integrals of the
motive $\bH^1(\bP^1\setminus\{0,1,\infty\})$, which are exactly the
motivic multiple zeta values (Brown 2012). Applying
$\Phi^{\mathrm{mot}}_*$ to the
$\widehat{\mathfrak{t}}_2$-presentation of $r_{\mathrm{bar}}(z)$
produces an expansion in $t_{12}^n$ with MZV coefficients at weight
$\geq 3$. The coefficient $\tau_{n_1,\ldots,n_r}$ is extracted by
composing the iterated OPE bracket of depth $r$ with the motivic
comodule action. The $\mathrm{GRT}^{\mathrm{fin}}$-projection kills
the MZV terms (they lie in the commutator subalgebra), so for
algebras where the face identity holds through the abelianization
(affine Kac--Moody, Heisenberg) the motivic face agrees with the
bar face; for class $M$ algebras (Virasoro, $\Walg_N$) where the
higher commutator terms are detected, $\mathrm{F}_8$ carries
genuinely new MZV data.
\end{proof}

\subsection{The operadic face $\mathrm{F}_9$}\label{sec:operadic-face}

\begin{theorem}[Operadic face]\label{thm:operadic-face}
Let $\mathrm{GC}_2$ be the Kontsevich graph complex of stable
$2$-dimensional graphs, and let
$\mathrm{GRT}_1(\bQ)\cong H^0(\mathrm{GC}_2)$ be Willwacher's
isomorphism (\emph{Grothendieck--Teichmüller group and
Kontsevich graph complex}, Invent.\ Math.\ 2015). For every graph
cocycle $\gamma\in\mathrm{GC}_2$, the corresponding
$\mathrm{GRT}$-element $\Phi^\gamma$ acts on the bar complex
$B(\cA)$ by a graph-summed derivation. The \emph{operadic face} is
\[
 \mathrm{F}_9(\cA)(z)
 \;=\;
 \sum_{\Gamma}\frac{1}{|\mathrm{Aut}\,\Gamma|}\,
 c_\Gamma(z)\cdot W_\Gamma(\cA),
\]
summed over stable graphs $\Gamma$ with two external legs, where
$c_\Gamma(z)$ is the graph weight (a rational function of $z$
computed from Kontsevich configuration integrals) and
$W_\Gamma(\cA)$ is the corresponding multi-OPE bracket on $\cA$.
The face $\mathrm{F}_9$ is obtained from $\mathrm{F}_1$ by the
$\mathrm{GRT}$-element dual to the graph cocycle.
\end{theorem}

\begin{remark}[Operadic interpretation]
Face $\mathrm{F}_9$ realizes $r(z)$ as a graph-sum over Feynman
diagrams in the Kontsevich formality quasi-isomorphism for the
$E_2$-operad. It is the face that sees the operadic (topological
configuration-space) origin of the collision residue most directly:
each stable graph contributes a universal rational function
$c_\Gamma(z)$, and the sum over graphs is controlled by the
graph-complex cohomology. In contrast to $\mathrm{F}_1$
(bar-intrinsic: coordinate-dependent but canonical) and
$\mathrm{F}_8$ (motivic: period-valued), $\mathrm{F}_9$ is
manifestly configuration-space-geometric.
\end{remark}

\subsection{The landscape census reconciliation (FM99)}
\label{sec:grt-landscape-fm99}

The Vol I landscape census lists the affine Kac--Moody collision
residue in two forms that appear contradictory:
\[
 r(z) \;=\; \frac{k\,\Omega}{z}
 \qquad\text{and}\qquad
 r(z) \;=\; \frac{\Omega}{(k+h^\vee)\,z}.
\]
These are \emph{not} a chained equality (as though
$k = 1/(k+h^\vee)$, which is false). They are \emph{two faces} in
the $\mathrm{GRT}^{\mathrm{fin}}$-torsor of
Theorem~\ref{thm:grt-face-torsor}(v), related by the Casimir
rescaling $\lambda = 1/(k(k+h^\vee))$:
\begin{equation}\label{eq:casimir-rescaling}
 \Phi^{\mathrm{KM}}_{k\to 1/(k+h^\vee)}\,\cdot\,
 \bigl(k\,\Omega\bigr)
 \;=\;
 \frac{\Omega}{k+h^\vee}.
\end{equation}
The left-hand side is the bar-intrinsic (unshifted) face produced
by the ordered bar propagator with propagator normalization
$1/z$; the right-hand side is the Sugawara-normalized face produced
by absorbing the level shift into the Casimir tensor. Both
present the same underlying element; the equality is a
\emph{tensor identity in the $\mathrm{GRT}^{\mathrm{fin}}$-quotient},
not a numerical identity of coefficients.

\begin{remark}[How this resolves FM99]
The failure mode FM99 recorded the attempted chain
$k\,\Omega/z = \Omega/((k+h^\vee)z)$ as a numerical identity and
found it false (at $k=1$, $h^\vee=2$ for $\fsl_2$:
$1\cdot\Omega = \Omega/3$ is false as scalars). The
reconstitution is: these are two orbit representatives, and the
$\mathrm{GRT}$-torsor is exactly the group that intertwines them.
AP126 (level-stripped $r$-matrix) is compatible with this
reconstitution: both faces satisfy $k=0\Rightarrow r=0$ (the
left face vanishes because its numerator is $k\,\Omega$; the
right face vanishes because the Sugawara normalization is
undefined when $k+h^\vee=\infty$ at $k=0$ via the reciprocal).
\end{remark}

\subsection{The $\mathrm{F}_7$ disambiguation and the super-variant}
\label{sec:f7-disambiguation}

Two distinct structures have been called ``face $\mathrm{F}_7$'' in
the literature and in earlier drafts (FM101):
\begin{itemize}
 \item $\mathrm{F}_7^{\mathrm{Gaudin}}$: the simple-pole projection
 $r_1$ (the generating function of the Gaudin Hamiltonians
 $H_i = \sum_{j\neq i}\Omega_{ij}/(z_i-z_j)$). This lives on the
 $E_\infty$ shadow and is defined for all four shadow classes.
 \item $\mathrm{F}_7^{M,\mathrm{top}}$: the top-degree $A_\infty$
 operation $m_3$ detecting non-formality (the class-$M$-specific
 face). This lives on the $E_1$ datum itself and is nonzero only
 for class $M$ (Virasoro, $\Walg_N$).
\end{itemize}

\begin{proposition}[$\mathrm{F}_7$ disambiguation]\label{prop:f7-disamb}
The Gaudin generator $\mathrm{F}_7^{\mathrm{Gaudin}}$ is the simple-pole
coefficient $r_1$ of the collision expansion
$r(z) = \sum_{k\geq 0}r_k\,z^{-k-1}$; it is
$\mathrm{GRT}^{\mathrm{fin}}$-invariant (a scalar per Casimir
rescaling) and is common to all four classes $G,L,C,M$. The
top-$A_\infty$ face $\mathrm{F}_7^{M,\mathrm{top}}$ is the class-$M$
extension: the operation $m_3\colon \cA^{\otimes 3}\to\cA$ that
encodes the non-formality of the bar complex at the highest
collision depth. $\mathrm{F}_7^{M,\mathrm{top}}$ is vacuous for
classes $G,L,C$ (where $m_3=0$ by formality) and carries genuine
new data for class $M$.
\end{proposition}

The naming is: $\mathrm{F}_7$ without superscript means
$\mathrm{F}_7^{\mathrm{Gaudin}}$ (the default); the top-$A_\infty$
face is always written $\mathrm{F}_7^{M,\mathrm{top}}$ or
equivalently $\mathrm{F}_7'$.

\subsection{The super-variant (AP107)}\label{sec:super-variant}

For chiral algebras with \emph{odd} generating fields (symplectic
fermions, $bc$ ghosts, super-affine Kac--Moody), the collision
residue from simple poles is distinct from the Laplace-transform
$r$-matrix (AP107). The bar-intrinsic construction is unchanged,
but the $\mathrm{F}_1\Leftrightarrow\mathrm{F}_2$ identification
requires a sign twist.

\begin{theorem}[Super-variant of
Theorem~\textup{\ref{thm:seven-faces}}]\label{thm:super-variant}
Let $\cA$ be a super-chirally-Koszul chiral superalgebra
(generating fields graded by $\bZ/2$). All statements of
Theorems~\textup{\ref{thm:seven-faces}} and
\textup{\ref{thm:grt-face-torsor}} hold verbatim after the following
substitutions:
\begin{enumerate}[label=\textup{(\alph*)}]
 \item The tensor product $\otimes$ is the $\bZ/2$-graded
 (Koszul-sign) tensor product.
 \item The Casimir $\Omega$ is the super-Casimir
 $\Omega^{\mathrm{super}} = \sum_a (-1)^{|a|}t_a\otimes t^a$,
 computed using the supertrace-invariant bilinear form.
 \item The associator gauge group is the super-$\mathrm{GRT}_1$
 acting on the parity-graded pure braid Lie superalgebra
 $\widehat{\mathfrak{t}}_n^{\mathrm{super}}$.
 \item The $\mathrm{F}_1\Leftrightarrow\mathrm{F}_2$ identification
 is mediated by the super-twist
 $\sigma_{\mathrm{super}}(a\otimes b) = (-1)^{|a||b|}b\otimes a$
 rather than the bosonic flip.
\end{enumerate}
The $k_{\max}=p_{\max}-1$ relation, the operator-order trichotomy,
and the four-class classification are unchanged. The
$\beta\gamma$ system and the symplectic fermion example are the
two archetypal class-$C$ super-landscape witnesses.
\end{theorem}

\subsection{Table of representatives}\label{tab:grt-representatives-sec}

\begin{center}
\renewcommand{\arraystretch}{1.35}
\begin{tabular}{|c|l|l|l|}
\hline
\textbf{Face} & \textbf{Associator} $\Phi$
 & \textbf{Gauge}
 & \textbf{Where this face is native} \\
\hline\hline
$\mathrm{F}_1$ & $1$ (trivial)
 & bar-intrinsic
 & universal (Vol~I Thm~A) \\
\hline
$\mathrm{F}_2$ & $\Phi^{\mathrm{DNP}}$
 & line-operator $R$-twist
 & 3d HT (DNP 2025) \\
\hline
$\mathrm{F}_3$ & $\Phi^{\mathrm{KZ-PVA}}$
 & $\lambda$-bracket
 & PVA (Khan--Zeng 2025) \\
\hline
$\mathrm{F}_4$ & $\Phi^{\mathrm{GZ26}}$
 & sphere Hamiltonian
 & holographic (GZ 2026) \\
\hline
$\mathrm{F}_5$ & $\Phi^{\mathrm{KZ}}$ (Knizhnik--Zamolodchikov)
 & Yangian $\hbar$
 & Yangian (Drinfeld 1985) \\
\hline
$\mathrm{F}_6$ & $\Phi^{\mathrm{STS}}$
 & Sklyanin flip
 & Poisson (STS 1983) \\
\hline
$\mathrm{F}_7$ & $\Phi^{\mathrm{FFR}}$
 & Casimir simple pole
 & Gaudin (FFR 1994) \\
\hline
$\mathrm{F}_7^{M,\mathrm{top}}$ & --- (not $\mathrm{GRT}$-conjugate)
 & class-$M$ $A_\infty$ top
 & Virasoro, $\Walg_N$ (this work) \\
\hline
$\mathrm{F}_8$ & $\Phi^{\mathrm{mot}}$ (Brown motivic)
 & MZV-valued
 & mixed Tate motives (Brown 2012) \\
\hline
$\mathrm{F}_9$ & $\Phi^{\mathrm{Will}}$ (Willwacher graph)
 & Kontsevich graph complex
 & $E_2$-formality (Willwacher 2015) \\
\hline
$\mathrm{F}_\Phi$ & any $\Phi\in\mathrm{GRT}_1(\bQ)$
 & torsor orbit
 & general associator gauge \\
\hline
\end{tabular}
\end{center}
\label{tab:grt-representatives}

The first seven rows are the historical faces of
Theorem~\ref{thm:seven-faces}. Rows 8--9 are the
$\mathrm{GRT}$-reconstitution faces of
Theorems~\ref{thm:motivic-face} and~\ref{thm:operadic-face}. The
bottom row is the general member of the $\mathrm{GRT}$-orbit; the
face family is infinite and parametrized by $\mathrm{GRT}_1(\bQ)$.
The row $\mathrm{F}_7^{M,\mathrm{top}}$ is \emph{not} in the
$\mathrm{GRT}$-orbit of $\mathrm{F}_1$: it is an independent face
detecting non-formality in class $M$ only, and the torsor
identification does not apply to it.

%%% ================================================================
%%% 4. THE TRICHOTOMY
%%% ================================================================

\section{The classification trichotomy}\label{sec:trichotomy}

\subsection{The four shadow classes}

The single-line dichotomy partitions the standard landscape into
four classes, classified by the shadow depth $r_{\max}$.

\begin{center}
\renewcommand{\arraystretch}{1.35}
\begin{tabular}{|c|c|c|l|}
\hline
\textbf{Class} & $r_{\max}$ & $\beta_{\mathrm{level}}$
 & \textbf{Examples} \\
\hline\hline
$G$ (Gaussian) & $2$ & quadratic
 & Heisenberg, lattice VOAs, free fermion, $bc$ ghosts \\
\hline
$L$ (Lie/tree) & $3$ & quadratic
 & Affine $\widehat{\fg}_k$ all simple types, non-critical \\
\hline
$C$ (Contact) & $4$ & simple-pole
 & $\beta\gamma$, symplectic fermions, matrix factorizations \\
\hline
$M$ (Mixed) & $\infty$ & higher-pole
 & Virasoro, $\Walg_N$, principal DS reductions \\
\hline
\end{tabular}
\end{center}

\noindent
The classes are mutually exclusive and exhaustive. Every
chirally Koszul algebra in the standard landscape falls into
exactly one class, and every class is realized by at least one
family.

\subsection{The three invariants table}

For each family in the standard landscape, the triple
$(p_{\max}, k_{\max}, r_{\max})$ takes an explicit value that
determines the class and the operator order of the GZ26
Hamiltonians.

\begin{center}
\renewcommand{\arraystretch}{1.35}
\begin{tabular}{|l|c|c|c|c|c|}
\hline
\textbf{Algebra} & $p_{\max}$ & $k_{\max}$ & $r_{\max}$
 & \textbf{Op.\ order} & \textbf{Class} \\
\hline\hline
Heisenberg $\cH_k$ & $2$ & $1$ & $2$ & $0$ & $G$ \\
\hline
Lattice $V_{\Lambda}$ & $2$ & $1$ & $2$ & $0$ & $G$ \\
\hline
Affine $\widehat{\fsl}_{N}{}_k$ & $2$ & $1$ & $3$ & $0$ & $L$ \\
\hline
$\beta\gamma$ & $1$ & $0$ & $4$ & trivial & $C$ \\
\hline
Virasoro $\mathrm{Vir}_c$ & $4$ & $3$ & $\infty$ & $2$ & $M$ \\
\hline
$\Walg_3$ & $6$ & $5$ & $\infty$ & $4$ & $M$ \\
\hline
$\Walg_N$ & $2N$ & $2N-1$ & $\infty$ & $2N-2$ & $M$ \\
\hline
Bershadsky--Polyakov $\Walg^{(2)}_3$
 & $4$ & $3$ & $\infty$ & $2$
 & $M$ ($T$) / $G$ ($J$) \\
\hline
\end{tabular}
\end{center}

\subsection{The operator-order trichotomy}

\begin{theorem}[Operator-order trichotomy]\label{thm:op-trichotomy}
For every $\cA$ in the standard landscape, the collision depth
$k_{\max}(\cA)\in\{0,1\}\cup\{3,4,5,\ldots\}$. The classification
into three operator regimes is:
\begin{enumerate}[label=\textup{(\roman*)}]
 \item \emph{Trivial}\/: $k_{\max}=0$. The GZ26 Hamiltonians are
 identically zero. Examples: $\beta\gamma$, symplectic
 fermions (class $C$).
 \item \emph{Multiplication}\/: $k_{\max}=1$. The GZ26
 Hamiltonians are multiplication operators of order $0$.
 Examples: Heisenberg (class $G$), affine Kac--Moody (class $L$).
 \item \emph{Differential}\/: $k_{\max}\geq 3$. The GZ26
 Hamiltonians are genuine differential operators of order
 $k_{\max}-1$. Examples: Virasoro (order $2$),
 $\Walg_N$ (order $2N-2$); both are class $M$.
\end{enumerate}
The value $k_{\max}=2$ does not occur in the standard landscape,
because it would require $p_{\max}=3$, which is forbidden by
bosonic locality and half-integer weight constraints.
\end{theorem}

\begin{proof}
The exclusion of $k_{\max}=2$ is Proposition~\ref{prop:no-k2}
below. The three cases are:
\begin{itemize}
 \item $k_{\max}=0$: the collision residue has only the
 constant ($z^0$) term, which by degree counting cannot produce
 a nontrivial differential operator at any order; therefore
 the Hamiltonians vanish.
 \item $k_{\max}=1$: the collision residue has a simple pole
 at $z=0$, which by residue extraction gives a multiplication
 operator. For affine Kac--Moody, this is the Casimir.
 \item $k_{\max}\geq 3$: the collision residue has poles of
 order $\geq 3$, whose residue extraction gives differential
 operators of positive order. For Virasoro at $k_{\max}=3$,
 the order is $2$; for $\Walg_N$ at $k_{\max}=2N-1$, the order
 is $2N-2$.
\end{itemize}
The match with the four-class shadow classification is direct:
classes $G$ and $L$ both have $k_{\max}=1$; class $C$ has
$k_{\max}=0$; class $M$ has $k_{\max}\geq 3$.
\end{proof}

\begin{proposition}[Absence of $k_{\max}=2$]\label{prop:no-k2}
The value $k_{\max}=2$ does not occur in the standard landscape.
\end{proposition}

\begin{proof}
The generator OPE pole order $p_{\max}$ takes values in
$\{1,2\}\cup\{4,5,\ldots\}$ for bosonic chiral algebras in the
standard landscape. The OPE of two weight-$h$ currents has poles
at orders $0,1,\ldots,2h$, and the maximal nonzero pole is $2h$
only for bosonic self-OPEs. For $h=1$ (currents), $p_{\max}=2$;
for $h=2$ (stress tensors), $p_{\max}=4$; intermediate values of
$p_{\max}=3$ would require half-integer weight, which breaks
bosonic locality. Since $k_{\max}=p_{\max}-1$, the corresponding
excluded value is $k_{\max}=2$.
\end{proof}

\subsection{Why $\beta\gamma$ is exceptional}

The $\beta\gamma$ system is the critical witness for the
independence of $r_{\max}$ from $p_{\max}$. Its shadow depth is
$r_{\max}=4$ (class $C$), but its generator OPE pole order is
$p_{\max}=1$ (not $2$, not $4$). The reason: the generators
$\beta$ and $\gamma$ have the simple-pole OPE
$\beta(z)\gamma(w)\sim (z-w)^{-1}$, with $\beta\beta$ and
$\gamma\gamma$ both regular. The Sugawara energy-momentum
tensor $T = :(1-\lambda)\beta\partial\gamma - \lambda\partial\beta\gamma:$
is a composite field of weight $2$ whose self-OPE has a quartic
pole $(z-w)^{-4}$, but $T$ is \emph{not} a generator: it is
determined by $\beta$ and $\gamma$ and does not appear in the
generating set.

The consequence for the seven-face programme is that the
GZ26 commuting Hamiltonians of the $\beta\gamma$ system vanish
identically: there is nothing to quantize at the sphere level,
because the collision residue $r(z)$ of $\Theta_{\beta\gamma}$ is
zero. The shadow depth $r_{\max}=4$ of $\beta\gamma$ comes not
from higher OPE poles but from the quartic contact invariant,
which lives on a charged stratum of the cyclic deformation
complex and exits the complex by rank-one rigidity (stratum
separation).

\begin{remark}[Stratum separation]\label{rem:stratum-separation}
Class $C$ is characterized by a mechanism distinct from classes
$G$, $L$, and $M$: the termination of the shadow obstruction tower
at degree~$4$ occurs not because the shadow metric $Q_L$ becomes
a perfect square (as in $G$ and $L$) nor because it remains
infinite (as in $M$), but because the quartic contact invariant
lives on a stratum whose self-bracket exits the complex. On the
primary line, $Q_L$ is the constant $(2\kappa)^2$, but the full
cyclic deformation complex has additional charged strata that
carry the contact invariant. This is why class $C$ appears in
the classification with the annotation $\Delta=0^{*}$: the
discriminant vanishes on the primary line but the contact
obstruction is nonzero in the full complex.
\end{remark}

%%% ================================================================
%%% 5. COMPUTE ENGINES
%%% ================================================================

\section{Compute engines}\label{sec:compute}

The seven-face programme is backed by fourteen compute engines,
each independently verifying one or more face identifications.
Every engine has a test file with multiple independent
verification paths per tested formula.

\begin{center}
\renewcommand{\arraystretch}{1.3}
\begin{tabular}{|l|c|l|}
\hline
\textbf{Engine} & \textbf{Tests} & \textbf{Faces} \\
\hline\hline
\texttt{theorem\_gz26\_commuting\_hamiltonians\_engine} & $32$ & F4, F7 \\
\hline
\texttt{theorem\_dnp\_meromorphic\_tensor\_engine} & $23$ & F1, F2 \\
\hline
\texttt{theorem\_bv\_brst\_genus1\_constraints\_engine} & $59$ & F3 (genus 1) \\
\hline
\texttt{theorem\_w3\_holographic\_datum\_engine} & $83$ & F4, F7 ($\Walg_3$) \\
\hline
\texttt{theorem\_pva\_classical\_r\_matrix\_engine} & $61$ & F3, F5 \\
\hline
\texttt{theorem\_virasoro\_spectral\_r\_matrix\_engine} & $91$ & F5 (Vir) \\
\hline
\texttt{theorem\_nonprincipal\_line\_operators\_engine} & $63$ & F2 (non-principal DS) \\
\hline
\texttt{theorem\_genus2\_planted\_forest\_gz26\_engine} & $73$ & F4 (genus 2) \\
\hline
\texttt{theorem\_w3\_commuting\_hamiltonians\_engine} & $105$ & F4, F7 ($\Walg_3$) \\
\hline
\texttt{theorem\_three\_way\_r\_matrix\_engine} & $72$ & F5, F6, F7 \\
\hline
\texttt{theorem\_ainfty\_nonformality\_class\_m\_engine} & $77$ & trichotomy (M) \\
\hline
\texttt{theorem\_dk0\_evaluation\_bridge\_engine} & $63$ & F2, F4 \\
\hline
\texttt{theorem\_sl3\_yangian\_r\_matrix\_engine} & $88$ & F5 ($\fsl_3$) \\
\hline
\texttt{theorem\_three\_paper\_intersection\_engine} & $39$
 & invariants, trichotomy \\
\hline\hline
\multicolumn{1}{|r|}{\textbf{Total}} & $\mathbf{929}$ &
 \multicolumn{1}{r|}{} \\
\hline
\end{tabular}
\end{center}

\noindent
The critical engine is
\texttt{theorem\_three\_paper\_intersection\_engine}: its $39$
tests check the $(p_{\max},k_{\max},r_{\max})$ triple for every row
of the three-invariants table and confirm the operator-order
assignment. The $\beta\gamma$ row is verified by the dedicated
test \texttt{test\_betagamma\_class\_C}, which is the sharpest
cross-check for the three-invariants distinction: it confirms
$p_{\max}(\beta\gamma)=1$, $k_{\max}(\beta\gamma)=0$,
$r_{\max}(\beta\gamma)=4$, operator order trivial, class $C$.

%%% ================================================================
%%% 6. STATUS AND LOOSE ENDS
%%% ================================================================

\section{Status and loose ends}\label{sec:status}

\subsection{Volume status}

The seven-face programme is proved for the standard landscape in
the three-volume monograph. The per-face per-volume status is:

\begin{center}
\renewcommand{\arraystretch}{1.3}
\begin{tabular}{|l|c|c|c|l|}
\hline
\textbf{Face} & \textbf{Vol I} & \textbf{Vol II} & \textbf{Vol III}
 & \textbf{Source} \\
\hline\hline
F1 (twisting morphism) & proved & specialized & specialized
 & Vol I Thm A \\
\hline
F2 (line operator R-twist) & cited & proved & specialized
 & DNP 2025 \\
\hline
F3 (PVA $r$-matrix) & proved & used & specialized
 & KZ 2025 \\
\hline
F4 (GZ26 commuting) & proved & used & specialized
 & GZ 2026 \\
\hline
F5 (Yangian $r$-matrix) & proved & cited & specialized
 & Drinfeld 1985 \\
\hline
F6 (Sklyanin bracket) & proved & cited & specialized
 & STS 1983 \\
\hline
F7 (Gaudin generator) & proved & cited & specialized
 & FFR 1994 \\
\hline
\end{tabular}
\end{center}

\subsection{Loose ends}

Three loose ends remain in the seven-face programme:

\begin{enumerate}
 \item \emph{Non-algebraic-family constructions}\/: the master
 theorem is proved for algebraic families with rational OPE
 coefficients. The residual universality conjecture
 restricts to non-algebraic-family constructions, and no known
 modular Koszul example with simple Lie symmetry falls outside
 the algebraic-family class.
 \item \emph{Multi-weight genus $g\geq 2$}\/: for multi-weight
 algebras at $g\geq 2$, the face identifications receive a
 cross-channel correction $\delta F_g^{\mathrm{cross}}\neq 0$.
 The seven-face programme is the uniform-weight-lane projection;
 the
 full multi-weight generalization requires an additional
 summand from cross-channel stable graphs and is the subject
 of the multi-weight genus expansion theorem.
 \item \emph{Non-principal $W$-algebras}\/: arbitrary-nilpotent
 BV duality (bar-cobar commutes with non-principal DS
 reduction) is conjectural beyond hook-type in type~$A$. The
 seven-face programme is proved for principal DS reductions
 and for the hook-type corridor in type $A$.
\end{enumerate}

The loose ends do not affect the master theorem on the standard
landscape; they precisely record the scope boundaries.

\subsection{Anti-patterns to avoid}

Three anti-patterns expose the common failure
modes in using the seven-face programme:

\begin{itemize}
 \item Never conflate the three invariants
 $p_{\max}$, $k_{\max}$, $r_{\max}$. Writing ``shadow depth''
 when one means ``OPE pole order'' is wrong in general: the
 $\beta\gamma$ system has shadow depth $4$ and OPE pole
 order $1$.
 \item Distinguish generator OPE pole order
 from composite OPE pole order. The Sugawara tensor of
 $\beta\gamma$ has quartic self-OPE, but the generators
 $\beta$ and $\gamma$ have simple-pole OPE. The bar complex
 is built from generators, not composites.
 \item Do not equate shadow depth with operator
 order without the $k_{\max}$ translation. The GZ26 operator
 order is $\max(k_{\max}-1,0)$, not $r_{\max}-2$. The two
 agree for classes $G$, $L$, $M$ but disagree for class $C$.
\end{itemize}

%%% ================================================================
%%% 7. CLOSING REMARKS
%%% ================================================================

\section{Closing remarks}\label{sec:closing}

The seven-face programme is the unifying thread of the three-volume
monograph on chiral algebras and modular Koszul duality. Its
content is not a new theorem in any single framework: each face
$\mathrm{F}i\Leftrightarrow\mathrm{F}j$ identification uses
established tools in its source framework. Its content is the
\emph{network} of identifications: the assertion that seven
different mathematical traditions, working in seven different
categorical languages with seven different computational
apparatuses, all describe the same underlying element
$r(z)\in\cA^{!}\otimes\cA^{!}[\![z^{-1}]\!]$ of the Koszul dual.

The surprising aspect is not that equivalences exist: algebraic
structures often have multiple presentations. The surprising
aspect is that the equivalences close the chain without
obstruction, so that starting from any one face and walking the
chain returns to the same underlying object. This is the
\emph{coherence} of the programme: every face is verified against
every other face, both in the mathematics of the monograph and in
the $929$-test compute layer. Coherence at this scale is not
automatic; it is earned, face by face, by checking that the
normalization conversions (classical $\hbar\leftrightarrow$
quantum $\hbar$, PVA $\lambda$-brackets $\leftrightarrow$ OPE
modes, Yangian level shift $\leftrightarrow$ affine level shift)
all agree under cross-composition.

The four-class classification $G$, $L$, $C$, $M$ gives the seven
faces a structural skeleton. Classes $G$ and $L$ are the
``easy'' cases: the shadow metric on the primary line is a
perfect square, the shadow obstruction tower terminates at degree~$3$, and the
GZ26 Hamiltonians are multiplication operators. Class $M$ is the
``generic'' case: the shadow obstruction tower is infinite, the Hamiltonians
are differential operators of order $2N-2$, and the seven faces
all carry nontrivial information. Class $C$ is the
``exceptional'' case: the shadow obstruction tower terminates at degree~$4$ by
stratum separation, but the collision residue vanishes, so the
GZ26 Hamiltonians are trivial. The $\beta\gamma$ system is the
archetype of class $C$, and it is also the critical witness for
the independence of the three invariants.

We close with the observation that the seven-face programme is an
\emph{audit instrument}: every time a new face (an eighth framework,
a new quantization scheme, a new categorical presentation) is
added to the network, its identification with the other seven must
be proved separately, and the compute engines must be extended to
verify it. The programme is not closed: it is open-ended, and
each new face is a new testable prediction. The current
statement of Theorem~\ref{thm:seven-faces} is the snapshot at
seven; future work will extend it to eight or more.

%%% ================================================================
%%% REFERENCES
%%% ================================================================

\begin{thebibliography}{99}

\bibitem{DNP25}
T.~Dimofte, N.~Niu, A.~Py, \emph{Line operators in $3$d
holomorphic-topological gauge theory and categorified Koszul
duality}, preprint 2025.

\bibitem{KZ25}
Z.~Khan, P.~Zeng, \emph{Poisson vertex algebras and $3$d sigma
models with holomorphic-topological boundary}, preprint 2025.

\bibitem{GZ26}
D.~Gaiotto, P.~Zeng, \emph{Holographic reconstruction of chiral
algebras from commuting sphere Hamiltonians}, preprint 2026.

\bibitem{Drinfeld1985}
V.~G.~Drinfeld, \emph{Hopf algebras and the quantum Yang-Baxter
equation}, Dokl.\ Akad.\ Nauk SSSR \textbf{283} (1985), 1060--1064.

\bibitem{STS1983}
M.~A.~Semenov-Tian-Shansky, \emph{What is a classical $r$-matrix?},
Funct.\ Anal.\ Appl.\ \textbf{17} (1983), 259--272.

\bibitem{FFR1994}
B.~Feigin, E.~Frenkel, N.~Reshetikhin, \emph{Gaudin model, Bethe
ansatz and critical level}, Comm.\ Math.\ Phys.\ \textbf{166}
(1994), 27--62.

\bibitem{Lorgat2026a}
R.~Lorgat, \emph{Modular Koszul duality: a monograph on bar-cobar
duality for chiral algebras on algebraic curves}, Volume~I
(manuscript, 2026).

\bibitem{Lorgat2026b}
R.~Lorgat, \emph{$A$-infinity chiral algebras and $3$d
holomorphic-topological QFT}, Volume~II (manuscript, 2026).

\bibitem{Lorgat2026c}
R.~Lorgat, \emph{Calabi--Yau quantum groups and the categorified
seven-face programme}, Volume~III (manuscript, 2026).

\end{thebibliography}

\end{document}
